\bookchapter{제39장: 컨텍스트 엔지니어링 — AI가 내 의도를 이해하게 만들기}

\chapterintro{이런 경험 해 보셨나요? 새 대화 창으로 바꾸니 AI가 기억을 잃은 것처럼 프로젝트의 기본 상황을 몽땅 잊어버리는 거예요. AI가 일부러 등을 쓰는 게 아니라, 원래 '건망증'이 있기 때문이에요. 이번 장에서는 AI에 '장기 기억 시스템'을 장착하는 법을 알려 주고, 아울러 2026년 가장 뜨거운 개발 패러다임 중 하나인 스펙 기반 개발도 만나 보게 될 거예요.}

\booksection{프로젝트 기억 파일: AI의 '장기 기억'}

AI 모델에는 근본적인 제약이 하나 있어요. 기억이 '단기'라는 거죠. 대화를 시작할 때마다 AI는 여러분의 프로젝트에 대해 알는 게 하나도 없어요. 그래서 대화 속에서 계속 '우리는 React를 써요', '백엔드는 Node.js예요', '데이터베이스는 PostgreSQL이에요', '코드 스타일은 큰따옴표예요'라고 알려 줘야 하죠. 이건 비효율적일 뿐 아니라 빠뜨리기 쉬워요. 새로 온 동료에게 매일 입사 교육을 다시 시키는 것과 같으니까요.

해결책은 '프로젝트 기억 파일'이에요. 프로젝트 루트 디렉터리에 두는 일반 텍스트 파일로, AI 도구가 대화를 시작할 때마다 자동으로 읽어요. 도구마다 이름 규칙은 조금씩 달라요.

\begin{center}
\begin{tabularx}{\linewidth}{|X|X|}
\hline
\textbf{도구} & \textbf{파일 이름} \tabularnewline
\hline
Cursor & \texttt{.\allowbreak cursorrules} \tabularnewline
Claude Code & \texttt{CLAUDE.\allowbreak md} \tabularnewline
GitHub Copilot & \texttt{copilot-\allowbreak instructions.\allowbreak md} \tabularnewline
범용 방안 & \texttt{AGENTS.\allowbreak md} \tabularnewline
\hline
\end{tabularx}
\end{center}

\begin{center}\small\color{gray}도구별 프로젝트 기억 파일 이름 규칙\end{center}

이름은 달라도 원리는 같아요. AI에게 '항상 기억해 주길' 바라는 정보를 이 파일에 적는 거예요. 좋은 프로젝트 기억 파일에는 이런 내용이 들어가야 해요. 프로젝트 개요(어떤 프로젝트인가), 기술 스택(무엇을 쓰는가), 코드 규범(명명, 들여쓰기, 주석 스타일), 디렉터리 구조(코드가 어떻게 조직되어 있는가), 개발 규칙(브랜치 전략, 커밋 규범, 테스트 요구). Karpathy의 CLAUDE.md는 겨우 70줄이지만 핵심 정보를 모두 담고 있어요. 간결하지만 완전하죠.

\begin{lstlisting}
# 프로젝트: 개인 가계부 앱
## 개요
사용자가 일상 수입·지출을 기록하도록 돕는 풀스택 개인 가계부 앱이에요.
## 기술 스택
- 프런트엔드: React + TypeScript + Tailwind CSS
- 백엔드: Node.js + Express + TypeScript
- 데이터베이스: SQLite(개발) / PostgreSQL(운영)
- 테스트: Vitest
## 코드 규칙
- 큰따옴표 사용
- 들여쓰기 2칸
- 함수 이름은 camelCase
- 타입 정의는 PascalCase
- 모든 함수에 JSDoc 주석 필수
## 디렉터리 구조
/src
  /frontend  - 프런트엔드 코드
  /backend   - 백엔드 코드
  /shared    - 프런트·백 공유 타입
  /tests     - 테스트 파일
## 개발 규칙
- 커밋 메시지는 한국어로
- 새 기능은 테스트를 먼저 작성
- API 응답은 { success, data, error } 형식으로 통일
\end{lstlisting}

\begin{tcbwarn}{주의: '두 번째 요구사항 문서'로 만들지 마세요}
이 파일의 길이에 주의하세요. 수십 페이지짜리 요구사항 문서로 쓰면 안 돼요. 프로젝트 기억 파일의 목적은 'AI가 배경을 빠르게 파악하게 하기'이지 '모든 기능을 자세히 묘사하기'가 아니에요. 간결하고 구조적으로 유지하고, 'AI가 알아야 할 것'만 적으세요.

\end{tcbwarn}

\booksection{스펙 기반 개발(SDD)}

Spec-Driven Development(SDD, 스펙 기반 개발)는 2026년에 점차 성숙해 가는 개발 패러다임이에요. 핵심 철학을 한 문장으로 말하면, 스펙을 먼저 쓰고 코드를 나중에 쓴다는 거예요. 전통적인 TDD(테스트 기반 개발)와 달리 SDD의 '스펙'은 테스트 케이스만이 아니라 요구사항 설명, 인터페이스 정의, 데이터 모델 등도 포함해요. AI는 이 스펙에 따라 코드를 생성하고, 사람의 역할은 스펙을 작성하고 유지하는 거죠.

SDD의 작업 흐름은 다섯 단계로 진행돼요.

\begin{enumerate}
  \item 자연어로 요구사항 스펙을 적어요(.md 파일 하나에 작성).
  \item 인터페이스와 데이터 구조를 정의해요(스펙 파일에 작성).
  \item AI가 스펙에 따라 코드 구현을 생성하게 해요.
  \item 테스트를 돌려 결과물이 스펙에 맞는지 검증해요.
  \item 맞지 않으면 스펙을 고치거나 AI에게 코드 수정을 시켜요.
\end{enumerate}

이 흐름에서 스펙 파일이 '유일한 진리의 원천'이에요. 코드는 다시 생성할 수 있지만, 스펙은 먼저 올바라야 하죠.

스펙 파일 작성

스펙이 충분한가?

충분함

AI가 스펙에 따라 코드 생성

테스트를 실행해 산출물 검증

미흡 / 테스트 실패

코드 문제 코드 수정; 스펙 문제 스펙 수정

검증으로 돌아가기

모두 통과 납품

\begin{center}\small\color{gray}그림 3: 스펙 기반 개발(SDD)의 작업 흐름\end{center}

SDD의 강점은 스펙 파일이 코드보다 안정적이라는 점이에요. 기반 모델이나 도구가 바뀌어도 여러분의 스펙 파일은 여전히 유효하고, 새로운 AI 도구에게 코드만 다시 생성시키면 돼요. 이 책이 '도구보다 방법론이 중요하다'고 강조하는 이유이기도 해요. 스펙 파일이야말로 여러분의 방법론을 구체적으로 담는 그릇이니까요.

\booksection{컨텍스트 윈도우 관리}

프로젝트 기억 파일이 있어도 대화 중에 '컨텍스트가 부족한' 문제를 만나게 돼요. AI 모델의 컨텍스트 윈도우(한 번에 처리할 수 있는 정보량)는 유한하거든요. 대화가 너무 길거나 참조하는 파일이 너무 많으면 AI가 앞의 정보를 '잊어버려서' 결과물 품질이 떨어져요. 친구와 6시간 동안 수다를 떨면 맨 처음에 했던 말을 친구가 까맣게 잊은 것과 같죠.

컨텍스트 윈도우 관리의 핵심 전략은 '계층화된 참조'예요.

\begin{center}
\begin{tabularx}{\linewidth}{|X|X|X|}
\hline
\textbf{컨텍스트 층위} & \textbf{내용} & \textbf{관리 방식} \tabularnewline
\hline
1층: 프로젝트 기억 & 기술 스택, 규범, 디렉터리 구조 & CLAUDE.md 등의 파일에 적어 두고 자동으로 불러오므로 반복할 필요 없음 \tabularnewline
2층: 파일 참조 & 현재 과제와 관련된 코드 파일 & @ 기호로 정확히 참조하고, 전문을 대화에 붙여 넣지 않음 \tabularnewline
3층: 대화 컨텍스트 & 현재 과제의 대화 기록 & 과제 하나가 끝나면 새 대화를 열어 컨텍스트가 비대해지는 걸 방지 \tabularnewline
4층: 즉시 지시 & 지금 이 단계의 구체적인 요구 & 명확하고 간결하며 목적에 맞게, 장황한 묘사는 피함 \tabularnewline
\hline
\end{tabularx}
\end{center}

\begin{center}\small\color{gray}컨텍스트 윈도우의 계층별 관리 전략\end{center}

\booksection{컨텍스트 엔지니어링 실전 포인트}

컨텍스트 엔지니어링은 한 번 하면 끝나는 설정이 아니라 계속 다듬어 가는 과정이에요. 다음은 실전에서 반복적으로 검증된 포인트들이에요.

\begin{enumerate}
  \item 기억 파일은 프로젝트가 발전함에 따라 함께 업데이트하세요. 새 기술을 도입하거나 아키텍처를 바꿀 때는 대화 속에서 임시로 설명하는 대신, 제일 먼저 기억 파일을 업데이트해요.
  \item @ 참조를 잘 활용하고 전문 붙여 넣기는 자제하세요. 대부분의 AI 프로그래밍 도구는 @file, @folder 문법을 지원하며, 백그라운드에서 파일 내용을 AI의 컨텍스트에 주입해요. 손으로 붙여 넣는 것보다 효율적이고 대화를 오염시키지도 않죠.
  \item 대화는 '짧고 굵게' 하세요. 한 대화에서는 과제 하나 또는 기능 포인트 하나만 다뤄요. 끝나면 새 대화를 열어 다음을 진행하고요. 긴 대화는 컨텍스트 윈도우를 소모할 뿐 아니라 AI의 '집중'도 흐트러뜨려요.
  \item 결과를 '설명'하지 말고 AI에게 '보여' 주세요. 버그를 만나면 '어떤 오류가 있는 것 같아요'라고 자기 말로 묘사하는 대신, 오류 로그를 @로 바로 참조해 주세요. AI는 여러분의 전달보다 원본 정보를 직접 보는 편이 훨씬 효과적이에요.
  \item 주기적으로 정리하고 보관하세요. 끝낸 대화와 임시 파일은 작업 공간을 차지해요. 관련 없는 내용을 주기적으로 치우고 프로젝트 디렉터리를 깔끔하게 유지하면 AI가 프로젝트 구조를 더 효율적으로 이해하는 데도 도움이 돼요.
\end{enumerate}

\begin{tcbpractice}{직접 해 보기: 내 프로젝트에 70줄 기억 파일 쓰기}
지금 하고 있거나(하려고 계획 중이거나) 하는 프로젝트 하나를 골라 CLAUDE.md를 써 보세요. 프로젝트 개요(2\textasciitilde{}3문장), 기술 스택(목록으로), 코드 규범(3\textasciitilde{}5개), 디렉터리 구조(트리로 그리기), 개발 규칙(2\textasciitilde{}3개). 70줄 안으로 유지하세요. 다 쓰고 나면 프로젝트에 대한 이해 자체가 훨씬 명확해진 자신을 발견하게 될 거예요.

\end{tcbpractice}

\begin{tcbsummary}{이번 장 정리}
\begin{itemize}
  \item AI의 기억은 '단기'예요: 새 대화가 시작되면 프로젝트에 대해 아는 게 0이에요
  \item 프로젝트 기억 파일(CLAUDE.md 등)은 AI의 '장기 기억'으로, 대화마다 자동으로 불러와요
  \item 좋은 기억 파일 = 개요 + 기술 스택 + 규범 + 디렉터리 + 규칙, 70줄이면 충분해요
  \item SDD 스펙 기반 개발: 스펙을 먼저 쓰고 코드를 나중에 써요. 스펙은 '유일한 진리의 원천'으로 코드보다 안정적이에요
  \item 컨텍스트는 4층으로 관리해요: 프로젝트 기억 @파일 참조 대화 기록 즉시 지시
  \item 실전 포인트 5가지: 함께 업데이트, @참조, 짧은 대화, 원본 전달, 주기적 정리
\end{itemize}

\end{tcbsummary}
